\documentclass[11pt,a4paper]{article}
\usepackage[margin=21mm]{geometry}
\usepackage{amsmath,amssymb,xcolor}
\definecolor{navy}{RGB}{24,49,70}
\setlength{\parindent}{0pt}
\setlength{\parskip}{5pt}
\makeatletter
\renewcommand\section{\@startsection{section}{1}{0pt}{10pt}{4pt}{\large\bfseries}}
\makeatother
\begin{document}
{\color{navy}\LARGE\bfseries D4: Keeping the locator correct}\\[5pt]
{\large Emergency Triage Queue --- Oracle Owner}

\section*{What are we proving?}
Structure A stores patients in a heap array $h$. The locator $pos$ tells us
where each patient is stored. For every valid index $i$, we need:
\[
\boxed{pos[h[i].id]=i}
\]
This means: if a patient is at index $i$, looking up their ID must return $i$.
We assume each live patient has a unique ID.

\section*{1. After a swap}
Suppose patient 42 is at index 2 and patient 90 is at index 6. When we swap
them, we also set $pos[42]=6$ and $pos[90]=2$.
Both mappings are now correct. Other patients did not move, so their mappings
are still correct too. The same argument works for any two indices.

\section*{2. After a sift}
Sift-up and sift-down are just sequences of swaps. The locator is correct
before the first swap, and each swap keeps it correct. Therefore, it remains
correct after the whole sift. During build-heap, we first create the mappings
for the input array, then use these same swaps.

\section*{3. After deleting a patient}
First, use the locator to find the patient. Swap that patient with the last
record and update both mappings. Then remove the last record and erase its
ID from the locator. If a replacement remains, sift it up or down as needed
to restore heap order.

The swap keeps the mappings correct, removal deletes only the target's
mapping, and the sift preserves the remaining mappings. If the patient was
already last, no sift is needed. In a linear-probing table, deletion leaves
a tombstone so searches can still reach other IDs in the same probe chain.

\section*{4. When is the indexed heap valid?}
We need both heap order and a correct locator:
\begin{itemize}
\item \textbf{Heap order:} every parent comes before or ties with its children
under the full comparison rule: priority key, then arrival time, then ID.
\item \textbf{Locator consistency:} each patient has a unique ID mapped to
its exact index, with no extra live mappings for absent patients.
\end{itemize}
These conditions are \textbf{necessary}: broken heap order makes the heap
invalid, while a wrong locator can send an ID lookup to the wrong record.
They are also \textbf{sufficient}: heap order makes the root a minimum
because it comes before every descendant, and the locator finds every
patient correctly. The array already provides the complete binary-tree shape.

\textbf{Main idea:} Whenever a patient moves, update their locator entry.
That is what keeps the heap and locator consistent.
\end{document}
